% Relative to the directory from which pdfLaTeX runs.
\makeatletter
\def\input@path{{../../../preamble/}{../../../../book2026-2027/preamble/}}
\makeatother
\input{preamble.tex}

% Add only the shared graph inputs used by this activity.
\geometry{margin=0.75in}
\newif\ifprintsolutions
\printsolutionsfalse % Use \printsolutionstrue to include the review pages.

\begin{document}

\begin{center}
    {\Large\bfseries [Activity title]\par}
    % \medskip [Learning criteria, when supplied]\par
\end{center}

% For a guided lesson, add the stages the lesson needs.
% \section*{Method}
% [Introduce the method or notation.]
% \section*{Worked Example}
% [Demonstrate the method and interpret the result.]
% Add partially scaffolded tasks before the less prompted practice tasks.

\section*{Problems}
\begin{enumerate}[leftmargin=*, itemsep=1em]
    % Use a minipage for a short problem. Remove it for a problem that may span pages.
    \item \begin{minipage}[t]{\linewidth}
        \textbf{[Problem title]}\par
        [Context, givens, units, and constraints.]

        \medskip
        \begin{tabularx}{\linewidth}{@{}>{\raggedright\arraybackslash}X
                                           >{\raggedright\arraybackslash}X@{}}
            \toprule
            \textbf{[Data heading and units]} & \textbf{[Student entry and units]} \\
            \midrule
            {[Given value]} & \rule{0pt}{2.4em} \\
            \bottomrule
        \end{tabularx}

        \begin{enumerate}[label=(\alph*), leftmargin=*]
            \item {[Task using the table.]}
            \item {[Task requiring a written response.]}\par
                \vspace{1.2cm}
                \noindent\rule{\linewidth}{0.4pt}
        \end{enumerate}
    \end{minipage}
\end{enumerate}

\ifprintsolutions
    % If the student pages use \newgeometry, use \restoregeometry here.
    \clearpage
    \section*{Solutions}
    \begin{enumerate}[leftmargin=*, itemsep=1em]
        \item \textbf{[Problem title]}
            \begin{enumerate}[label=(\alph*), leftmargin=*]
                \item {[Connect the givens to the method and show the working.]}
                \item {[Interpret the result in context.]}
            \end{enumerate}
    \end{enumerate}
\fi

\end{document}
